\diarytitle{0063}{熱伝導方程式の初期値・境界値問題（ディリクレ境界条件、2 モードの重ね合わせ）}

境界条件がディリクレ型（$u(0,t)=u(L,t)=0$）なので、変数分離から固有値問題までの手順は次の通りである。$u(x,t)=X(x)T(t)$ とおき方程式に代入すると
\[
\frac{T'(t)}{k\,T(t)} = \frac{X''(x)}{X(x)} = -\lambda
\]
となり、$X''+\lambda X=0$、$X(0)=X(L)=0$ の固有値問題を解くと、$\lambda \le 0$ では自明解しか得られず、$\lambda_n=(n\pi/L)^2$、$X_n(x)=\sin(n\pi x/L)$（$n=1,2,3,\dots$）が固有値・固有関数である。時間成分は $T_n(t)=e^{-k\lambda_n t}$。よって一般解は
\[
u(x,t) = \sum_{n=1}^\infty c_n \sin\frac{n\pi x}{L}\,e^{-k(n\pi/L)^2 t}
\]
である。$t = 0$ とおいて初期条件と比べると
\[
\sum_{n=1}^\infty c_n \sin\frac{n\pi x}{L} = 3\sin\frac{\pi x}{L} - 2\sin\frac{3\pi x}{L}
\]
右辺はすでに固有関数 $\sin\dfrac{\pi x}{L}$（$n=1$）と $\sin\dfrac{3\pi x}{L}$（$n=3$）の線形結合なので、フーリエ係数の積分計算をするまでもなく係数比較で
\[
c_1 = 3, \qquad c_3 = -2, \qquad c_n = 0\ (n \neq 1, 3)
\]
と定まる（固有関数系の直交性により、この表し方は一意）。各モードは自分の固有値に応じた速さ $e^{-k(n\pi/L)^2 t}$ で独立に減衰するから、それぞれの係数に指数因子を付ければよい。よって
\[
\boxed{\; u(x,t) = 3\sin\frac{\pi x}{L}\,\exp\!\left(-k\left(\frac{\pi}{L}\right)^{2} t\right) - 2\sin\frac{3\pi x}{L}\,\exp\!\left(-k\left(\frac{3\pi}{L}\right)^{2} t\right) \;}
\]

（検算: 各項は問1と同じ形の固有モードであり、線形性から $u_t = k u_{xx}$ を満たす。境界: $x=0,L$ で $\sin(n\pi x/L)=0$ なので $u(0,t)=u(L,t)=0$。初期条件: $t=0$ で $e^{0}=1$ となり $u(x,0)=3\sin\frac{\pi x}{L}-2\sin\frac{3\pi x}{L}$ に一致。）
